\section*{开始前：这本讲义怎样读}
\addcontentsline{toc}{section}{开始前：符号、证明和学习顺序}
这份讲义按老师原课件的次序重写。正文不需要另开原课件：定义、结论、证明和纠错都写在这里。相邻续页合成一个完整论证；各节标题标原页范围，末尾还有逐页索引。\textbf{“原页”是教师 PDF 的页码，不是旧 gpt 笔记的页码。}本讲义自己的页码在页脚。

“思路”说明为什么选择这条路线；“证明”是需要能重写的内容；“批注”解释选点、抽子列或条件用途；“补充”表示教师没有展开但理解所必需的论证。原稿有错误时明确给出修正，不能为了照抄路线而保留错误等式。

建议第一遍连读定义与例子，第二遍把证明的假设、目标、关键构造单独抄出来，第三遍盖住答案做最后的练习。暂时不熟悉的基础可以先用下面的工具，不要在主证明中反复猜符号。

\subsection*{1．点、集合、对应：先看对象类型}
$x\in A$ 表示一个点属于一个集合；$A\subseteq B$ 表示 $A$ 中\emph{每个}点都属于 $B$。证明 $A=B$，通常先证 $A\subseteq B$ 再证 $B\subseteq A$。$A\setminus B$ 是在 $A$ 中而不在 $B$ 中的点，不是数的减法。

函数 $f:S\to T$ 每个输入给一个输出点；对应 $\Phi:S\rightrightarrows T$ 每个输入给一个输出集合 $\Phi(x)\subseteq T$。$y\in\Phi(x)$ 表示 $y$ 是输入 $x$ 下的一种允许选择。$\Phi(A)=\bigcup_{x\in A}\Phi(x)$，表示从输入集合 $A$ 中任取输入时能得到的全部输出。图是输入输出对的集合：
\[G(\Phi)=\{(x,y):x\in S,\ y\in\Phi(x)\}.\]
注意 $\Phi(x)$ 在输出空间中，$G(\Phi)$ 在输入与输出的乘积空间中，闭值与闭图是两回事。

\subsection*{2．量词如何指导证明}
“$\forall x\ \exists y$”表示给定任意 $x$ 后可以选一个依赖于它的 $y$，不保证同一个 $y$ 对所有 $x$ 都有效。证明全称命题时写“任取”；证明存在命题时给构造或调用保证存在的定理。
\[\neg(\forall x\ P(x))\iff\exists x\ \neg P(x),\qquad
\neg(\exists x\ P(x))\iff\forall x\ \neg P(x).\]
例如“每个邻域都遇到 $E$”的否定是“存在一个邻域完全不遇到 $E$”。不要把否定后仍写成“每个”。

要反证“附近全部值集都在开集 $G$ 内”，否定式是：任意小的输入邻域里都能找一个坏输入，并在它的值集中找一个跑出 $G$ 的点。于是对半径 $1/n$ 分别选
\[\|x_n-x_0\|<1/n,\quad y_n\in\Phi(x_n)\setminus G.\]
这就是证明里突然出现两列点的原因：它们是把否定式逐个应用在半径 $1,1/2,1/3,\ldots$ 上。

\subsection*{3．开、闭、紧分别给你什么}
开集给\textbf{空间余量}：$x\in G$ 时存在 $r>0$ 使开球 $N(x,r)\subseteq G$。若 $x_n\to x\in G$，收敛定义保证充分大的 $n$ 有 $x_n\in G$。

闭集给\textbf{极限归属}：若 $x_n\in F$ 且 $x_n\to x$，则 $x\in F$。\textbf{闭性不保证点列收敛。}例如闭集 $\R$ 中的列 $n$ 没有收敛子列。

紧集给\textbf{有限覆盖和子列}：每个开覆盖可选有限子覆盖；每个集内序列都有收敛子列，且极限在该集内。有限维欧氏空间中，紧等价于闭且有界。证明中“有界”先给 B--W 子列，随后“闭”才使极限留在集合里。

子列写作 $x_{n_k}$，其中 $n_1<n_2<\cdots$。如果原列收敛，每个子列收敛到同一极限；\textbf{一个子列收敛不能推出原列收敛}。有限多个坐标要一起收敛，必须逐次在已有子列里再抽：先使第一坐标收敛，再在这条子列中使第二坐标收敛。这样先前的收敛不会丢失。

\subsection*{4．两个常用补步：距离逼近、三角不等式}
定义 $\dist(y,A)=\inf_{a\in A}\|y-a\|$。$A\ne\varnothing$ 时，对每个 $\eps>0$ 都存在 $a\in A$ 使
\[\|y-a\|<\dist(y,A)+\eps.\]
否则所有距离都至少为右边那个更大的数，与下确界是最大下界矛盾。证明常用 $\eps=1/n$；这不要求存在最近点。

当 $z_n$ 接近 $y_n$、$y_n$ 又接近 $y$ 时，用
\[\|z_n-y\|\le\|z_n-y_n\|+\|y_n-y\|\]
将两段误差加起来。两项各小于 $\eps/2$ 就保证总误差小于 $\eps$。这解释为什么老师总把精度分成一半或三分之一。

\subsection*{5．实值函数半连续，与集合值对应不要混用}
实值 $f$ 在 $x_0$ 下半连续指
\[\forall\eps>0\ \exists U\ni x_0\ \forall x\in U:\quad f(x)>f(x_0)-\eps.\]
上半连续把它改成 $f(x)<f(x_0)+\eps$。前者防止附近函数值明显低于旧值，后者防止明显高于旧值。
对于每条 $x_n\to x_0$，对应的序列判据分别为
\[f(x_0)\le\liminf_n f(x_n),\qquad
\limsup_n f(x_n)\le f(x_0).\]
这里不假定 $f(x_n)$ 已有极限。$\liminf$ 看尾部下界的极限，$\limsup$ 看尾部上界的极限；最优化证明中可完全使用上面的 $\eps$ 邻域语言。

非空紧集上的上半连续实值函数达到最大值，下半连续实值函数达到最小值。以最大值为例：上半连续性给每点附近一个有限上界，紧性把这些邻域有限化，于是全局有上界；再选越来越接近上确界的点，紧性抽子列，半连续性使极限点的函数值至少是上确界，所以取到最大值。每一步分别用了什么条件，应当能说出来。
